% Tags.tex — Documentation of % tags used in .nw source files
%
% These tags appear as %-comments in the .nw literate programming sources.
% They serve two purposes:
%   1. Provenance: who wrote or requested a paragraph
%   2. Theme: what kind of content the paragraph contains
%
% A paragraph can carry multiple tags on the same line, e.g.:
%   %claude: %cs-history:
% The first tag is always the provenance tag; theme tags follow.
%
% Future: theme tags may map to LaTeX environments (tcolorbox) so
% readers can visually identify (or skip) certain kinds of content.
% See the \ifwanthistorycode, \ifwantalternatives flags in Macros.tex.

%==========================================================================
% PROVENANCE TAGS — who wrote this paragraph
%==========================================================================
%
% Tag            Meaning
% -------------- ----------------------------------------------------------
% %claude:       Paragraph written by Claude (AI-generated prose).
%                REQUIRED on every AI-written paragraph.
%                Stays even when a theme tag is added alongside it.
%
% %yoann:        Author wants a %claude: paragraph added after this line.
% %yoann1/2/3:   Author wants a fix to the %claude: paragraph above.
%                (Delete the %yoann1/2/3: line after addressing it.)

%==========================================================================
% AUTHOR'S PRIVATE ANNOTATIONS — never remove or move
%==========================================================================
%
% Tag       Meaning
% --------- --------------------------------------------------------------
% \t        "thinking" or "todo" — author's reasoning, unfinished thoughts
% \l        "less important" — tangential asides, cross-references, lighter todo
% \n        "note" — topic tags, internal reminders
% %         Plain %-comments — author's own remarks

%==========================================================================
% THEME TAGS — what kind of content
%==========================================================================
%
% These classify the *topic* of a paragraph (written by claude or the author).
% They appear after %claude: on the same line if there is already a provenance
% tag for the paragraph. A paragraph should carry at most one
% theme tag (pick the dominant theme).
%
%--------------------------------------------------------------------------
% CORE TAGS — for LaTeX environments (skip/seek navigation)
%--------------------------------------------------------------------------
%
% These five tags are the highest priority. They will become visual
% LaTeX boxes (tcolorbox or similar) so that readers can quickly
% identify and either skip or focus on certain content.
%
% Paragraphs that DON'T carry any of these tags are the core
% content: Plan 9's specific approach and the code itself.
% That's what nobody skips.
%
% Tag              Skip if...                 Seek if...
% -----------      -------------------------  ----------------------------
% %terminology:    "I already know OS/        "I need to understand this
%                   compiler concepts"          term before reading code"
%
% %cs-history:     "I just want to            "I want the full
%                   understand the code"        intellectual context"
%
% %modern:         "I only care about         "I want to know how this
%                   Plan 9"                    relates to my daily tools"
%
% %others:         "I don't need the ARM      "I want to appreciate the
%                   vs x86 vs RISC-V           design space and Plan 9's
%                   details"                   choices"
%
% %plan9-is-cleaner: "I already know Plan 9   "I want to understand what
%                     is elegant, just show    Plan 9 redesigned with
%                     me the code"             decades of Unix hindsight"
%
%--------------------------------------------------------------------------
% Core Tag              Meaning                    Example
%--------------------------------------------------------------------------
%
% %cs-history:     Historical origin.         "gprof (Graham, Kessler,    [TODO]
%                  Who invented it, when,      McKusick, Berkeley, 1982)
%                  where, and why it           was the first widely-deployed
%                  mattered at the time.       profiler..."
%
% %modern:         Modern counterpart.        "GCC moved from yacc to     [TODO]
%                  How the wider world does    hand-written recursive
%                  it today; competition to    descent in 2004; Clang was
%                  the Plan 9 approach.        hand-written from the start."
%
% %others:         Other systems.             "Linux has spin_lock,       [TODO]
%                  How alternative systems     FreeBSD has mtx, Windows
%                  implement the same          has KSPIN_LOCK. ARM puts
%                  concept. Not necessarily    args in R0-R3, x86-64 uses
%                  modern — just different.    six registers, Windows x64
%                  Helps appreciate the        uses four different ones."
%                  design space and Plan 9's
%                  specific choices.
%
% %plan9-is-cleaner: Plan 9 improvement.     "Unlike Unix's errno (an     [TODO]
%                  Where Plan 9 redesigned    opaque number), Plan 9 uses
%                  something with decades     human-readable error strings
%                  of Unix hindsight. The     via errstr()." / "pread/pwrite
%                  Plan 9 authors knew what   take an explicit offset, making
%                  went wrong in Unix and     seek+read atomic — avoiding the
%                  fixed it.                  race that Unix's separate
%                  (Replaces older %unix:     seek+read allows."
%                  tag in Kernel.nw.)
%
%--------------------------------------------------------------------------
% Other Tag              Meaning                    Example
%--------------------------------------------------------------------------
%
% %wib:            Worse-is-better.           "5c has neither IR nor SSA.  [TODO]
%                  Plan 9's deliberate         It goes straight from AST to
%                  simplification vs. the      assembly — a deliberate
%                  mainstream approach.         trade-off..."
%                  Named after Richard
%                  Gabriel's 1989 essay.
%
% %evolution:      Lineage / evolution.       "Coroutines (Conway 1958) -> [TODO]
%                  How a concept changed       threads (1990s) -> async/await
%                  across decades; the arc     (C# 2012) -> algebraic
%                  from invention to today.    effects (OCaml 5, 2022)."
%
% %why-study:      Why study this.            "The programming model has not [DONE]
%                  Justification for           changed since 1985. Registers,
%                  studying an old or          load/store, condition flags,
%                  simplified approach.        branches..."
%
% %terminology:    Concept / terminology.     "A terminal is a device (or  [TODO]
%                  Introduces or clarifies     emulator); a shell is a
%                  a key concept, definition,  program that reads commands."
%                  or term the reader needs    / "A context switch is the
%                  before reading the code.    act of pausing one process
%                  (Core tag — see above.)     and resuming another."
%
% %cross-book:     Cross-book connection.     "The widget library solves,  [TODO]
%                  Linking concepts across     in miniature, the same
%                  different books in the      problems an OS solves:
%                  Principia series.           multiplexing a resource..."
%
% %comeback:       Technology comeback.        "Coroutines (Conway 1958)   [DONE]
%                  A concept that was          were displaced by kernel
%                  invented, fell out of       threads in the 1980s-90s,
%                  favor, then returned        then returned as goroutines,
%                  in a new form. The key      async/await, and algebraic
%                  pattern is: invented ->     effects — vindication of the
%                  displaced -> reborn.        original idea in modern dress."
%
% %why-win:        Why X won.                 "Ethernet won not because it [DONE]
%                  Explains why a specific     was the best LAN technology
%                  technology prevailed        — token ring had better
%                  over its competitors.       determinism — but because it
%                  Focus is on the reasons     was cheap, simple to wire,
%                  (economic, social,          and good enough at every
%                  technical, or accidental)   speed bump from 10 Mbit to
%                  not just the outcome.       100 Gbit."
%
% %road-not-taken: Unfulfilled promise.       "Segmentation offered         [TODO]
%                  An idea with genuine        fine-grained protection and
%                  merit that history          logical grouping — but paging
%                  passed over. Not a          won on simplicity, and the
%                  failed technology but       potential of hardware-enforced
%                  a path still waiting        module boundaries remains
%                  to be explored. An          unrealized. Maybe you'll be
%                  invitation to the reader.   the one to revive it."
%
% %reframe:        Unexpected analogy.        "A kernel can be read as     [TODO]
%                  Reframes something         an unusual concurrent
%                  familiar through an        program whose 'threads' are
%                  unexpected lens that       the per-process kernel
%                  opens the reader's mind.   stacks." / "A widget library
%                  "Read X as Y."             solves, in miniature, the
%                                             same problems an OS solves."
%
% %design:         Design principle.          "Each tiny language is       [TODO]
%                  A general software          deliberately incomplete.
%                  design insight, often       awk does not sort; sort does
%                  arising from the code       not search — they compose."
%                  but transcending it.

%--------------------------------------------------------------------------
% RELATIONSHIP TO EXISTING TAGS
%--------------------------------------------------------------------------
%
% %history:   Author's short telegraphic notes about history (e.g.,
%             "%history: unix V1"). NOT a theme tag. %cs-history: is
%             the rendered-prose counterpart for real paragraphs.
%
% %alt:       Author's short notes about alternatives (e.g.,
%             "%alt: GUI widget toolkit"). NOT a theme tag. %modern:
%             is the rendered-prose counterpart.
%
% %trans:     Transitional text between sections. Stays as-is.
% %toc:       Table-of-contents / overview content. Stays as-is.
% %dup:       Duplicated from another book. Stays as-is.
% %old:       Notes about original code before modifications. Stays as-is.
% %ocaml:     OCaml comparison notes. Stays as-is.
% %dead:      Dead code notes. Stays as-is.

%==========================================================================
% FUTURE: LaTeX ENVIRONMENTS
%==========================================================================
%
% When we are ready, each core tag maps to a tcolorbox environment:
%
%   \usepackage{tcolorbox}
%   \newtcolorbox{terminologybox}{
%     colback=green!3, colframe=green!40,
%     title=Concept, fonttitle=\bfseries\small
%   }
%   \newtcolorbox{cshistorybox}{
%     colback=gray!5, colframe=gray!50,
%     title=Historical Note, fonttitle=\bfseries\small
%   }
%   \newtcolorbox{modernbox}{
%     colback=blue!3, colframe=blue!40,
%     title=Modern Counterpart, fonttitle=\bfseries\small
%   }
%   \newtcolorbox{othersbox}{
%     colback=yellow!3, colframe=yellow!40,
%     title=Other Systems, fonttitle=\bfseries\small
%   }
%
% The boxes make it easy for a reader to skip (or seek out) certain
% kinds of content. The \ifwanthistorycode and \ifwantalternatives
% flags in Macros.tex could conditionally hide entire boxes.
%
% The syncweb/noweblatexpad toolchain would need a small extension
% to translate %claude: %cs-history: into \begin{cshistorybox}...
%
%==========================================================================
% FUTURE: \concept{} macro
%==========================================================================
%
% Currently .nw files use {concept} (curly braces) or {\em concept}
% or even sometimes \emph{concept} for new terminology.
% A semantic \concept{} macro would:
%   - Render as \emph{} now (backward compatible)
%   - Later: auto-generate index entries, glossary links,
%     marginal notes, or tooltips
%   - Work with %terminology: tagged paragraphs to build a
%     concept map of each book

%==========================================================================
% TAGGING PROCESS — lessons learned
%==========================================================================
%
% This section documents the process for systematically tagging
% paragraphs, so that the same approach can be followed on any machine.
%
% WHERE TO TAG
% ------------
% Focus on Principles sections (both the main "\section{XXX principles}"
% and per-chapter "\section{YYY principles}"). These are where
% cs-history, modern, others, terminology, and plan9-is-cleaner
% paragraphs cluster. Use this grep to find them all:
%   grep -rn 'section{.*principles' --include='*.nw'
%
% TAG PLACEMENT
% -------------
% Tags go on their own %-comment line, never mixed with prose text.
% Grouping tags on one line is fine: %claude: %cs-history:
%
% When the whole paragraph is about the tagged topic, put the tag
% on the %claude: line:
%   %claude: %cs-history:
%   Virtual memory was invented for the Atlas computer...
%
% When only part of a paragraph is about the tagged topic, insert
% the tag before the relevant sentence (not at the paragraph start):
%   ...the core idea is the same everywhere.
%   %others:
%   Different architectures spell the modes differently...
%
% This matters because tags will eventually become LaTeX environments
% (\begin{cshistorybox}...\end{cshistorybox}), so the tagged text
% should be a clean chunk that could plausibly be boxed.
%
% If the cross-system comparison or historical fact is just a brief
% inline parenthetical like "(Linux's struct file, BSD's vnode,
% Windows' HANDLE)", don't tag it — it's too small to be a box.
%
% WHAT NOT TO TAG
% ---------------
% Don't tag core content — paragraphs that explain what something IS,
% how it works mechanically, or what Plan 9 specifically does. These
% are the backbone the reader never skips.
%
% The %terminology: tag is reserved for paragraphs whose whole point
% is disambiguating confusing terms (shell vs terminal, interrupt vs
% exception vs trap, thread-safe vs reentrant, elapsed vs wall-clock
% time). Don't use it for every paragraph that introduces a concept —
% that would tag everything. Use \concept{} inline markup for the
% actual defined terms instead.
%
% DISTINGUISHING THE CORE TAGS
% ----------------------------
% %cs-history: — "invented by X at Y in year Z" — who/when/where
% %modern: — "today, the mainstream does it THIS way instead" —
%   current competition/replacement for Plan 9's approach
% %others: — "Linux calls it X, FreeBSD calls it Y, Windows calls
%   it Z" — same concept, different implementation across systems.
%   Not necessarily modern — just different. Includes cross-arch
%   comparisons (ARM vs x86 vs RISC-V).
% %plan9-is-cleaner: — "unlike Unix which does X, Plan 9 does Y
%   because..." — specific improvements Plan 9 made with hindsight.
%   Different from core content because the point is the comparison
%   with Unix, not the Plan 9 design on its own terms.
%
% The test: if you removed the paragraph, would the reader still
% understand the Plan 9 code? If yes, it's taggable (history,
% modern, others, plan9-is-cleaner). If no, it's core content.
%
% PROCESS FOR A NEW BOOK
% ----------------------
% 1. Find all principles sections: grep -n 'section{.*principles'
% 2. Read each section start to end
% 3. For each %claude: paragraph, ask: is this cs-history, modern,
%    others, plan9-is-cleaner, terminology, toc, or trans?
%    If none — it's core content, leave it untagged.
% 4. Place the tag at the right granularity (paragraph or sentence).
% 5. Commit per-book.
%
% BOOKS COMPLETED (systematic principles pass)
% ---------------------------------------------
% Kernel, Libcore, Shell, Compiler, Assembler, Linker, Build system,
% Debugger, Profiler, Graphics, Network, Utilities, Windows, Widgets.
% (Machine and Editors have no taggable content in principles.)
%
% TODO FILES
% ----------
% docs/latex/TODO-comeback.txt     — ~25 comeback arcs to write
% docs/latex/TODO-why-win.txt      — ~35 why-X-won stories to write
% docs/latex/TODO-why-study.txt    — ~30 justifications to write
% docs/latex/TODO-road-not-taken.txt — ~25 unfulfilled promises
% docs/latex/TODO-reframe.txt      — ~40 unexpected analogies
% docs/latex/TODO-cs-history.txt   — systematic pass still incomplete
